\section{Promoted Ag: when does dilute nickel add useful ethylene oxide output?}
\label{sec:silver}

\textbf{Summary.} Promoted Ag epoxidation catalysts must favor ethylene oxide (EO) over combustion at an economic rate, and they lose performance in service. Ni-dependent selectivity retention on a fully promoted catalyst is already reported: Kemp's 1995 Ni-containing Ag/Cs/Re/S catalyst lost 4.3 selectivity points over 50 days of accelerated aging versus 6.1 for its Ni-free parent, though the pair also differed in Cs content and post-treatment.\citep{kemp1995} Jalil and colleagues showed that dilute Ni suppresses combustion more than EO formation on supported Ag; their main materials lack the full Cs/Re promoter system.\citep{jalil2025} Chloride operation could equally explain a retention difference: measured and modeled chloride responses and prior lower-chloride disclosures show that it alone can change selectivity and output.\citep{iyer2021,santos2024,lee2024underchloride,modifier2024} The original contribution is therefore narrow: a composition- and processing-controlled, Ni-dependent relationship that predicts when changing chloride operation can substitute for Ni, or when a useful Ni advantage survives it; better controls alone do not establish it, and a practical benefit needs no unique microscopic explanation. The program compares ordinary drift of Ni-treated, sham-processed and untreated material over a fixed horizon, including the best tested chloride policy for each (Stages A--B), tests any difference against a calibrated chloride-response baseline (Stage C), and, if the loss is partly recoverable, predicts when a chloride recovery step pays off (Stage D). Success would deliver a reproducible formulation/policy comparison and a validated prediction of where any Ni advantage ends. As an illustration of the stakes, not a predicted improvement: if EO and complete combustion are the only carbon outlets, raising ethylene selectivity from 0.90 to 0.91 at fixed EO output lowers reaction-level ethylene use by 1.10\%, oxygen use by 4.40\% and \COtwo{} formation by 10.99\%.

\subsection{Prior work and the remaining question}

The Summary's calculation uses $1/S$ moles of ethylene, $0.5+3(1/S-1)$ moles of oxygen and $2(1/S-1)$ moles of \COtwo{} per mole of EO at fractional ethylene selectivity $S$. It excludes recycle, separation, utilities and catalyst manufacture, whose costs a measured improvement must also survive.

Kemp added Ni after preparation of the Ag/Cs/Re/S catalyst; lithium-containing precursors were also used. Selectivity was measured at 40\% oxygen conversion after reoptimizing. Both catalysts started within a reported 86.0--86.5\% range, and the bulk Ni:Ag atomic ratio was about 1:835. A solvent-treated, Cs-matched Ni-free control started worse and was not aged. The result is a difference between losses from separate starting values, not an exactly measured final-selectivity difference or a demonstrated lifetime extension, and it neither isolates the Ni effect nor reports useful integrated output.\citep{kemp1995}

In Jalil's work, organochloride raised selectivity further. Processing controls, model surfaces and Ni spectroscopy support a Ni-associated effect, but nominal loading does not count working surface Ni. In the supplementary chloride-response trace, the EO rate on NiAg first rises and then declines substantially: an evolving Ni--chloride interaction, not a stationary, durable output advantage on the proposed formulation.\citep{jalil2025,jalil2025si}

Hwang and colleagues link the distinctive Cs/Re/chloride response to promoter architecture and propose that otherwise unselective oxygen becomes an additional epoxidizing intermediate. This motivates asking whether Ni is compatible with that architecture, but is not direct evidence that Ni supplies oxygen to Re or preserves their contact.\citep{hwang2026}

Chloride operation could equally explain a retention difference. Adjusting retained chlorine (deposited from alkyl-chloride modifiers, removed by competing processes including alkanes such as ethane) can trade selectivity against useful EO rate; the optimum depends on material and conditions, and no universal coverage or modifier setting follows.\citep{iyer2021} Iyer and Bhan already predicted reactor output from independently measured oxidation, chlorine-exchange and support-mediated EO degradation kinetics, and Santos and colleagues showed that temperature, oxygen, water and \COtwo{} change the chloride response.\citep{iyer2021,santos2024} Prior disclosures describe lower-chloride aging policies (28-day tests, then extrapolation) and materials with a more stable optimum chloride requirement that hold work rate by adjusting temperature. Neither a crossover extrapolated well beyond the 28-day data nor a stable chloride optimum proves stable activity.\citep{lee2024underchloride,modifier2024}

Because selectivity alone can reward an almost inactive catalyst, the endpoint is sustained EO output. The program does not assume that Ni preserves a particular oxygen species or Ag--Re arrangement, and a mechanistic claim needs a discriminator beyond a selectivity curve.

\subsection{Stage A: validate the materials and define a fair output comparison}

\paragraph{Materials.}
Use an existing, reproducible Cs/Re-promoted Ag/$\alpha$-Al\textsubscript{2}O\textsubscript{3} reference and its validated epoxidation reactor; access to both is an entry requirement. Hwang's Ag/support-area dependence favors keeping this working architecture over moving the promoters onto a dilute colloidal-Ag support.\citep{hwang2026} From each parent preparation, make untreated, sham-processed and Ni-treated materials. The first Ni route is Kemp's post-impregnation with a nickel-nitrate/amine-containing solution and staged drying, adapted after a small compatibility check. Start near a bulk Ni:Ag atomic ratio of 1:800, an illustrative choice close to the historical loading, rather than assuming that Jalil's 1:200 optimum transfers. The sham receives every processing step except the Ni precursor; document any unavoidable solution-composition difference. Measure Ag, Ni, Cs, Re and other promoter contents and Ag size distributions; matched bulk contents do not establish identical promoter placement.

The sham is the causal comparator for the Ni-addition step. The untreated reference detects processing damage and represents ordinary operation of the starting catalyst; because a processing penalty can appear after the fresh check, it also runs over the full horizon. If the sham loses the reference's useful response, allow one correction for an identified preparation problem, then stop this route if the problem remains. A Ni/sham retention difference without a useful advantage over untreated operation is processing dependent, not an improvement of the starting catalyst. The Kemp route does not establish isolated Ni in Ag; comparing unpromoted Ag and Ni--Ag made the same way is needed only to attribute the result to the Ni function Jalil reported on fresh catalysts.

\paragraph{Measurement and operating conditions.}
Before assessing small selectivity differences, commission temperature uniformity, transport sensitivity and feed-to-analyzer response. Close the carbon balance against a calibrated flow tracer, counting ethane fed to remove chlorine as a carbon source and separating net formation from outlet concentration when \COtwo{} or EO is cofed. Report selectivity from consumed ethylene, because $r_{\rm EO}/(r_{\rm EO}+r_{\rm CO_2}/2)$ holds only under its carbon-source and product assumptions; unresolved cofeed contributions limit interpretation.

Condition all three materials by one protocol rather than forcing equal initial selectivity by different treatments. Map a small overlapping bracket of chloride settings with returns to a reference setting. Hold temperature, pressure, ethylene, oxygen, ethane and flow common for Ni and sham, and find a useful untreated setting within the same operating and tuning limits. Specify absolute pressure and the actual chloride feed, not only a nominal chloriding index. Set dwell times from observed rates and gas holdup, up to a fixed maximum; if reference returns remain history dependent, report age-specific trajectories rather than a composite steady curve. These data answer two questions: does Ni change EO and combustion rates at common conditions, and which material gives more accepted EO at its own chosen setting? A shifted chloride optimum can answer the second without altered oxygen chemistry. If both materials pass validation, no difference between the fresh materials does not stop the program, since retention may differ later.

\paragraph{Accepted output and cost boundary.}
Before the ordinary-operation comparison, fix the output requirements (ethylene-conversion interval, selectivity floor, minimum EO work rate, and temperature and feed limits) from the reproduced reference and the intended operating decision, not from whichever endpoint favors Ni, and fix a finite horizon $H$ that starts with initial conditioning. For policy $p$,
\begin{equation}
 P_p(H)=\frac{Q_p(H)}{m_{{\rm Ag},0}H},\qquad
 Q_p(H)=\int_0^H \dot n_{{\rm EO},p}^{\rm accepted}(t)\,dt.
 \label{eq:ag-productivity}
\end{equation}
EO is accepted only if it meets the product criteria, judged in windows compatible with analyzer response; off-spec product and downtime count as zero. Conversion, selectivity and work-rate requirements apply in production windows and a return-assessment window, both fixed in advance; temperature and feed limits apply throughout. Planned conditioning, recovery and return intervals may produce little or nothing but have fixed maximum durations and a fixed return deadline. A policy that fails a production window or misses the deadline is inadmissible, and a low-output interval cannot be relabeled afterward as planned recovery. The clock includes chloride searches that recur in the policy. Also report unfiltered net EO, failure times, temperature history, and output per total catalyst mass and bed volume, since Ag normalization can hide dilution.

For each admissible policy $p$ against a named tested comparator $b$, the preparation-cost boundary per bed campaign over $H$, at matched initial Ag inventory and with actual bed masses and volumes reported, is
\begin{equation}
 C_{{\rm prep},p}-C_{{\rm prep},b}
 \leq v_{\rm EO}\,[Q_p(H)-Q_b(H)]
 -\sum_k v_k\,[U_{k,p}(H)-U_{k,b}(H)]
 -v_L[L_p(H)-L_b(H)],
 \label{eq:ag-preparation-cost}
\end{equation}
where $v_{\rm EO}$ is the value per mole of accepted EO, $U_k$ are resource or disposal quantities (measured feed and modifier, utilities on a stated basis, valued emissions or waste) with unit costs $v_k$, and $L$ is operating handling labor at $v_L$ per labor-hour. Preparation cost is not counted again as an operating resource, and downtime, already charged through output, is not automatically labor. Without justified valuations, report the symbolic boundary and physical differences, not a monetary saving.

\subsection{Stage B: compare ordinary drift and operating policies}

Within each preparation block, assign independent charges to the Ni/sham common-setting pair and to the useful policies from Stage A, including at least one independently replicated untreated policy, all starting with the same conditioning and running for $H$. A common-setting arm may double as a policy arm only if its entire history is identical. The common-setting pair measures the effect of Ni addition under that exposure; the policy arms compare what an operator could use, and the two are not interchangeable.

Fix the untreated policy before aging from its own chloride response: one lower-chloride setting if it meets the requirements, otherwise a credible reference setting or scheduled adjustment, not a fitted percentage of another catalyst's optimum. Ni and sham get the same operating limits and tuning effort. If holding accepted conversion needs a different flow or temperature, that change and its consequences belong to the policy. The untreated policy is required even if its fresh response matches the sham's; it needs no full chloride or recovery factorial, but the best tested untreated policy appears explicitly in the rankings.

Compare cumulative EO and reactant use over $H$, not endpoint selectivity: a lower-chloride policy may pay a continuing selectivity penalty, not a one-time startup cost, and a temperature controller can hide activity loss unless its trajectory is kept. Higher selectivity with lower work rate or higher temperature remains a tradeoff. No accelerated exposure is needed. If neither drift nor a useful material or policy difference is resolved within $H$, end the program at that preparation and operating scope; a severe chloride pulse or an extrapolated crossover cannot substitute.

If an important loss develops, compare one mild recovery step with continuation on matched-age beds of both materials that share the common-setting history. Take the recovery chloride setting from the established bracket, hold other feeds fixed, then return both beds to the common reference setting for an observation window fixed in advance. Follow the recovery and continuation arms separately afterward, since treatment may change later drift. This tests recoverable conditioning under that intervention, not general regeneration. A failed recovery shows persistence over the tested duration, not irreversibility.

\subsection{Stage C: test whether the difference exceeds the known chloride response}

The baseline is each material's measured chloride/product response, informed by established chlorine-exchange and reactor models and calibrated on the actual formulation and feed range, including response times and, where conversion requires it, axial product gradients (less combustion makes less water and \COtwo{} in the bed, so equal inlet products need not mean equal local states).\citep{iyer2021,santos2024} An uncalibrated literature model is not a fair baseline.

Compare the EO--\COtwo{} trajectory during recovery with independently measured steady or matched-age chloride-response curves. A trajectory that follows the curves can still show a useful Ni-dependent recovery speed, since the curves do not set how fast the bed moves along them. Departure from a whole-bed curve can reflect axial mixtures of chlorine states, temperature or product gradients, so it does not identify promoter relocation. Add reproducible rate memory (a slowly changing state that makes the current rate depend on past chloride history, not a count of selective oxygen sites) only if the calibrated baseline misses by more than its prediction uncertainty and an important difference remains.

Use at most one further discriminator, chosen by the result. If less self-generated \COtwo{} seems to explain the Ni benefit, compare the pair under a common, externally controlled product burden and predict an untested burden, accounting for bed gradients. If the loss survives reconditioning and an inventory change is plausible, examine retained elements and Ag morphology on sister samples; unchanged total Re does not exclude redistribution, and a changed mean particle size need not explain all lost function. EO cofeed measures only net product sensitivity, and contact-time curvature cannot identify a secondary combustion flux.

A working-state assay is justified only if it separates the remaining alternatives and can detect the minority population invoked; an unchanged element-average spectrum cannot exclude it when its abundance is unconstrained. Chlorine stripping gives an operational inventory, and trifluoroethanol titration is destructive and estimator dependent: Iyer and Bhan treat total uptake as an upper-bound site estimate because inactive adsorption is possible.\citep{iyer2023} Neither gives an exact active-site count on Ni-containing material without new validation. Oxygen-isotope return measurements and a Ni-by-Re composition series are optional extensions. Isotope return needs a justified link between labeled and returning oxygen, a complete balance and a resolvable signal; neither high selectivity nor undetected gas-phase return proves Ni-to-Re oxygen transfer.

\subsection{Stage D: predict when a chloride recovery step pays off}

If drift is partly recoverable, fit, for Ni and sham, the matched-age recovery and continuation traces with the simplest time-dependent description consistent with both EO and \COtwo{} rates. A single relaxation time is allowed only if its approach and return are resolved and repeatable. Check it on a second training duration, then choose a withheld duration within the tested range. The prediction must not rely on rates identifying surface coverage.

Over a common remaining horizon $H_r$, the recovery benefit relative to continuation is
\begin{equation}
 B_m(t_R;H_r)=\int_0^{H_r}
 \left[\dot n_{{\rm EO},m,R}^{\rm accepted}(t;t_R)
 -\dot n_{{\rm EO},m,C}^{\rm accepted}(t)\right]dt,
 \label{eq:ag-recovery-benefit}
\end{equation}
where $m$ is Ni or sham, $t_R$ is the recovery duration, and $t$ is measured from the intervention at operating age $a$ (whereas $t$ in Eq.~\eqref{eq:ag-productivity} starts at conditioning), so $H_r=H-a$. The integral includes recovery and return time, not just recovered peak activity. Resources are compared over the same interval. The model predicts the sign of $B_m$ and the EO and \COtwo{} trajectories for both materials.

$B_m$ compares policies within one material. To rank across materials, convert it to absolute output,
\begin{equation}
 Q_{m,R}(H)=Q_{m,C}(H)+B_m(t_R;H-a),
 \qquad m\in\{\mathrm{Ni},\mathrm{sham}\},
 \label{eq:ag-absolute-recovery}
\end{equation}
using that material's continuation baseline with identical pre-intervention history, clock and inventory basis. Rank \emph{all admissible tested} policies (Ni ordinary and recovery, sham ordinary and recovery, untreated ordinary) at the same $H$ under the same acceptance and resource rules, with inventory-normalized outputs and the cost boundary on the matched-inventory basis. A larger within-material benefit does not decide absolute output. A comparison against ordinary Ni alone supports only that narrower conclusion; the ranking covers the tested policies, not an unmeasured optimum.

Recovery learned under the common-setting history cannot be assumed to transfer to another policy history; that would need new validation. The withheld test repeats recovery and matched continuation of both materials on independent preparation blocks, with return setting and acceptance windows fixed in advance, and checks both absolute outputs and the ranking against the ordinary policies, including untreated. Do not retune chloride afterward. Sham recovery beating ordinary Ni does not remove a Ni advantage if Ni recovery is better, so the Ni recovery policy stays in the ranking.

Success requires an observed benefit or limit within $H_r$ that changes the operating decision, predicted accurately enough to resolve the smallest important difference. A calibrated chloride-response model that already predicts the result supports an operating choice but ends the claim of an additional Ni-specific mechanism; a Ni-dependent rate-memory term must beat it on withheld prediction and identify a useful preparation or operating dependence. A periodic-service model needs repeated returns to the same productive state and measured recurrence costs, not one successful recovery.

If the ordinary difference is persistent rather than recoverable, replace this stage with one product burden or exposure duration, fixed in advance, that the remaining explanation predicts will keep or remove the Ni advantage; do not pursue the recoverable and persistent branches in parallel. A result that only reproduces historical Ni retention, disappears under chloride adjustment without a further important relationship, or needs an open-ended synthesis campaign should remain a bounded negative or confirmatory study and not be expanded.
